% Sections VI–VII — Discussion and Limitations, Conclusion
% Draft 2026-10-20.

\section{Discussion and Limitations}
\label{sec:discussion}

\textbf{Where the graph helps (RQ1, RQ2).} Graph retrieval is better where
the question names entities whose relations are recorded in ORKG. A comparison
query directly returns contributions that contain both methods. The retrieved values
keep their predicate, so a method a paper proposes is distinguishable from one
it uses as a baseline, and benchmark traversal brings in the reported scores. Dense retrieval has no
access to this structure. A paper's card lists many values, and its similarity
to a question about two methods is dominated by the topic, so dense retrieval
tends to return papers about only one of the methods. We attribute the
collapse of the weaker embedder on comparisons (0.166) to this. Conversely,
dense retrieval is better on broad topics and on finding a paper from words of
its title. In these cases text similarity is enough, while graph retrieval
depends on the surface form under which an entity is stored in ORKG and on
the LLM extracting it correctly.

\textbf{Complementarity.} The two systems share fewer than one of eight
papers and succeed on different question types, which is why their fusion
outperforms both. We kept the system itself graph-only, since every retrieved
paper can then be traced to the query path that found it. A hybrid is the
natural next step, but we evaluated it only at the retrieval level.

\textbf{Retrieval differences shrink in the answers.} Better evidence does not
always become a better answer. On claim questions graph retrieval supplies far
more deciding papers, but verdict accuracy does not improve, because the
models often decide from prior knowledge. Answer quality thus also depends
on whether the model uses the evidence it is given. The answer metrics are
also noisier than the retrieval metrics. With almost identical context, a
repeated run moved the mean strict score by 0.05 and reversed the sign of the
graph--dense difference. Comparisons of RAG configurations based on a single
run per system and differences of a few points should therefore be treated
with caution. Repeated runs and paired intervals are cheap with local models.

\textbf{The LLM judge (RQ3).} A judge that sees only source titles cannot
measure what distinguishes the retrievers, namely which papers are cited and
whether they support the answer. It tracked the gold metrics only moderately
and ranked models largely by whether they were cloud or local models. The
same judge model applied to claims with the full context, and required to
quote the context, gave usable per-model estimates. The evaluation protocol matters at
least as much as the judge model.

\textbf{Practical implications.} With good retrieval, a local 7B model came
within 0.07 of the best cloud model. The whole pipeline, including both
judges, runs locally, which matters when cloud free tiers change or withdraw
models.

\textbf{Limitations.}
\begin{itemize}
\item The question set is small (50 questions, 25 in the test split, 8--9 per
  type), so per-type results rest on few questions.
\item Relevance grades were proposed by an LLM and reviewed by the author. The
  blind check ($\kappa = 0.78$) was done by an annotator who had seen part of
  the grades before.
\item Run-to-run variance was measured for one model only, and the hybrid only
  at the retrieval level.
\item The results concern one graph, whose coverage and curation vary across
  fields. Cloud models were used through a free tier whose model offer changes.
\end{itemize}

\section{Conclusion}
\label{sec:conclusion}

In a controlled comparison over ORKG, graph-native and dense retrieval reached
the same overall quality but succeeded on different questions. The graph was
better for structured questions, embeddings for topical and title lookups, and
their fusion was best. At the answer level the difference was within the
variance between runs, and an LLM judge that does not see the evidence could
not measure it. Graded relevance gold sets and claim-level checks against the
context are a feasible alternative, and the evaluation should report repeated
runs.
